\rtmtight
\begin{tblr}{width=\textwidth,colspec={|Q[c,m,2.1cm]|X[l,m]|},hlines,vlines,hspan=minimal,rowsep=1.5pt,colsep=3pt,row{1}={bg=headgray,font=\bfseries}}
\SetCell{c,m}\hfil\logo\hfil & \textbf{POLITEKNIK NEGERI MALANG}\par\textbf{JURUSAN TEKNOLOGI INFORMASI}\par\textbf{PROGRAM STUDI: D4 TEKNIK INFORMATIKA} \\
\SetCell[c=2]{l} \textbf{RENCANA TUGAS MAHASISWA (RTM) DAN RUBRIK PENILAIAN} & \\
Nama Mata Kuliah & Internet of Things \\
Kode Mata Kuliah & RTI256105 \quad \textbf{sks / Jam perminggu:} 3 SKS/6 jam \quad \textbf{Semester:} 6 \\
Dosen Pengampu & Tim Pengajar Program Studi D4 TI \\
\end{tblr}
\assessmentblock{Tugas Perancangan Arsitektur IoT}{Case Method dan praktik perancangan}{SCPMK0204-04901, SCPMK0503-04902}{Mahasiswa merancang arsitektur sistem IoT end-to-end dan mengimplementasikan komponen komunikasi IoT berdasarkan studi kasus yang diberikan.}{Menganalisis kebutuhan sistem IoT, merancang arsitektur layer edge-fog-cloud, mengimplementasikan protokol komunikasi, menguji konektivitas, dan mendokumentasikan.}{Dokumen arsitektur IoT, diagram jaringan, kode implementasi protokol, dan laporan pengujian.}{Ketepatan rancangan arsitektur IoT end-to-end & 11 & Arsitektur mencakup layer sensor, gateway, fog, dan cloud dengan spesifikasi komponen, protokol, dan alur data yang tepat dan lengkap. & Arsitektur tersedia tetapi beberapa layer, komponen, atau spesifikasi protokol masih kurang detail atau kurang tepat. & Arsitektur tidak mencakup semua layer IoT atau tidak dapat menjadi acuan implementasi. \\ \\
Implementasi protokol komunikasi IoT & 11 & Protokol (MQTT/HTTP/CoAP) diimplementasikan dengan benar, koneksi broker stabil, topic/endpoint terstruktur, dan data terkirim sesuai spesifikasi. & Implementasi protokol berjalan tetapi beberapa konfigurasi, struktur topic, atau penanganan error belum optimal. & Protokol tidak dapat terkoneksi atau data tidak terkirim sesuai spesifikasi. \\ \\
Kualitas dokumentasi dan analisis kasus & 6 & Dokumen arsitektur lengkap, diagram jaringan informatif, dan analisis kasus menunjukkan pemahaman kebutuhan pengguna IoT. & Dokumentasi tersedia tetapi beberapa komponen diagram atau analisis kasus masih kurang rinci. & Dokumentasi tidak mencerminkan arsitektur yang dirancang atau analisis kasus tidak relevan. \\}

\assessmentblock{Kuis Konsep dan Protokol IoT}{Kuis}{SCPMK0204-04901, SCPMK0503-04902}{Kuis tertulis untuk mengukur pemahaman konsep IoT, arsitektur, dan protokol komunikasi yang diterapkan dalam sistem IoT.}{Menjawab soal pilihan ganda dan/atau uraian singkat terkait konsep, arsitektur, dan protokol IoT.}{Jawaban tertulis kuis.}{Pemahaman konsep dan arsitektur IoT & 5 & Konsep IoT, layer arsitektur, dan karakteristik komponen dijelaskan dengan benar dan tepat sesuai konteks soal. & Sebagian besar konsep dipahami tetapi penerapan pada arsitektur atau komponen masih kurang akurat. & Konsep IoT tidak dipahami atau jawaban tidak relevan dengan arsitektur IoT. \\ \\
Penguasaan protokol komunikasi IoT & 5 & Perbedaan, kelebihan, dan penggunaan tepat protokol MQTT, HTTP, CoAP, dan WebSocket dijelaskan secara akurat. & Sebagian protokol dipahami tetapi perbandingan atau konteks penggunaan masih kurang tepat. & Protokol komunikasi IoT tidak dapat dibedakan atau penjelasan keliru secara fundamental. \\}

\assessmentblock{UTS Perancangan Sistem IoT}{UTS berbasis perancangan}{SCPMK0204-04901, SCPMK0503-04902}{Evaluasi tengah semester untuk mengukur kemampuan merancang dan menganalisis arsitektur sistem IoT lengkap beserta implementasi protokol komunikasi.}{Merancang arsitektur sistem IoT sesuai studi kasus UTS, mengimplementasikan komponen komunikasi, dan mendokumentasikan keputusan teknis.}{Dokumen perancangan arsitektur IoT, diagram sistem, dan implementasi protokol komunikasi.}{Kelengkapan rancangan arsitektur sistem IoT & 8 & Rancangan mencakup semua layer IoT, spesifikasi komponen lengkap, dan alur data dari sensor ke cloud terdokumentasi dengan jelas. & Rancangan mencakup komponen utama tetapi beberapa layer atau spesifikasi masih kurang lengkap. & Rancangan tidak mencakup semua layer atau tidak dapat dijadikan acuan implementasi. \\ \\
Ketepatan pemilihan dan konfigurasi protokol & 7 & Protokol dipilih sesuai kebutuhan (latensi, bandwidth, QoS), konfigurasi broker/endpoint benar, dan justifikasi pemilihan tepat. & Protokol dipilih dengan alasan yang masuk akal tetapi konfigurasi atau justifikasi teknis masih kurang kuat. & Protokol dipilih tanpa justifikasi teknis atau konfigurasi tidak sesuai kebutuhan sistem IoT. \\}

\assessmentblock{PBL Implementasi Prototipe Sistem IoT}{Project Based Learning}{SCPMK0503-04902, SCPMK0904-04903}{Mahasiswa mengimplementasikan prototipe sistem IoT fungsional dari sensor hingga cloud menggunakan platform IoT modern dan mengintegrasikan layanan monitoring data real-time.}{Merakit rangkaian sensor/aktuator, memprogram mikrokontroler, mengkonfigurasi platform IoT, mengintegrasikan cloud, menguji end-to-end, dan mendokumentasikan.}{Prototipe IoT berjalan, kode sumber, dashboard monitoring, laporan implementasi, dan video demo.}{Fungsionalitas prototipe end-to-end & 10 & Prototipe berjalan dari sensor ke cloud, data terbaca akurat, kontrol aktuator responsif, dan alur data real-time terpantau. & Prototipe sebagian berjalan tetapi beberapa komponen, akurasi data, atau responsivitas kontrol masih bermasalah. & Prototipe tidak dapat mengirimkan data sensor ke cloud atau kontrol aktuator tidak berfungsi. \\ \\
Konfigurasi platform IoT dan integrasi cloud & 10 & Platform IoT (Node-RED/AWS IoT) dikonfigurasi lengkap, alur data divisualisasikan, alert terset, dan integrasi cloud berfungsi stabil. & Platform IoT terkonfigurasi dan berjalan tetapi beberapa fitur visualisasi, alert, atau integrasi cloud belum optimal. & Platform IoT tidak terkonfigurasi dengan benar atau integrasi cloud tidak berjalan. \\ \\
Dokumentasi dan kualitas implementasi & 5 & Kode terdokumentasi, wiring diagram tersedia, laporan implementasi lengkap, dan video demo menunjukkan fungsionalitas penuh. & Dokumentasi tersedia tetapi beberapa komponen seperti wiring diagram, kode, atau laporan masih kurang. & Dokumentasi tidak mencukupi untuk mereproduksi prototipe atau video demo tidak menunjukkan fungsionalitas. \\}

\assessmentblock{UAS Demo Sistem IoT}{UAS berbasis demo dan presentasi}{SCPMK0904-04903}{Mahasiswa mendemokan sistem IoT lengkap yang sudah dikonfigurasi dan diintegrasikan dengan cloud, mempresentasikan arsitektur, analisis data, dan kemampuan monitoring real-time.}{Demo sistem IoT live, presentasi arsitektur dan konfigurasi platform, menjawab pertanyaan teknis dari penilai.}{Demo sistem IoT live, bahan presentasi, dan laporan deployment.}{Kualitas demo sistem IoT end-to-end & 9 & Demo menampilkan aliran data real-time dari sensor ke dashboard cloud, kontrol aktuator berfungsi, dan sistem stabil selama demo. & Demo berjalan tetapi ada ketidakstabilan data, delay berlebihan, atau beberapa fitur tidak dapat didemonstrasikan. & Demo tidak dapat menampilkan aliran data real-time atau sistem tidak stabil saat dipresentasikan. \\ \\
Konfigurasi dan keamanan platform IoT & 7 & Platform dikonfigurasi dengan praktik keamanan yang baik (autentikasi, enkripsi), fitur lengkap, dan dapat di-maintain. & Platform berjalan tetapi konfigurasi keamanan atau kelengkapan fitur masih kurang optimal. & Platform tidak mengimplementasikan keamanan dasar atau konfigurasi tidak dapat di-maintain. \\ \\
Kemampuan analisis dan presentasi teknis & 6 & Arsitektur dan keputusan teknis dijelaskan dengan tepat, analisis data sensor bermakna, dan pertanyaan teknis dijawab secara kompeten. & Presentasi cukup jelas tetapi analisis data atau penjelasan beberapa keputusan teknis masih kurang mendalam. & Tidak dapat menjelaskan arsitektur sistem atau menjawab pertanyaan teknis dasar IoT. \\}

\begin{tblr}{width=\textwidth,colspec={|Q[l,m,3.2cm]|X[l,m]|},hlines,vlines,hspan=minimal,row{1-3}={bg=headgray},rowsep=1.5pt,colsep=3pt}
Jadwal Pelaksanaan: & Minggu 1--17 sesuai RPS. Setiap evaluasi memiliki RTM dan rubrik multi-kriteria. \\
Lain-Lain Yang Diperlukan: & LMS, repositori/kode sumber, perangkat lunak sesuai mata kuliah, dan perangkat presentasi. \\
Pustaka: & Mengacu pustaka utama dan pendukung pada bagian RPS. \\
\end{tblr}
